\documentclass[11pt,a4paper]{article}
\usepackage{templates/modeling-notes}
\usepackage{needspace}
\hypersetup{pdftitle={Sheet 1 -- Primal simplex}}
\title{\textbf{Mathematical Modeling}\\[2mm]
\large Sheet 1: Primal Simplex, Alternative Optima, and Growth Rays}
\author{Personal study notes}
\date{}
\begin{document}
\maketitle
{\small\tableofcontents}

\section{The task and the meaning of a simplex step}\label{sec:task}
\begin{otherlanguage}{polish}
  Rozwiąż podstawowym algorytmem sympleks. O ile to możliwe wystartuj z punktu
  $x_0$ Czy są rozwiązania alternatywne? W przypadku nieograniczonej funkcji
  podaj promień wzrostu.
\end{otherlanguage}
All six problems have three original decision variables. They are continuous
real quantities; no physical units are supplied. Every stated nonnegativity
condition is part of feasibility. For inequalities of the form $a^Tx\leq b$,
write $a^Tx+s=b$ with $s=b-a^Tx\geq0$. For $a^Tx\geq b$, write
$a^Tx-s=b$ with $s=a^Tx-b\geq0$. The different sign records which residual
is nonnegative. Auxiliary variables have objective coefficient zero and are
numbered after $x_3$ separately in each problem.

A basis $B$ selects three independent columns of the equality matrix $A$.
Set all nonbasic variables to zero and solve $A_Bx_B=b$. This produces a
basic feasible solution only when $x_B\geq0$. A feasible point need not be
basic, and an inadmissible supplied point cannot start primal simplex.
The basis order is the order of the rows in each displayed tableau.

\subsection{Why the signs and the ratio test work}\label{sec:step}
Multiplying the equations by $A_B^{-1}$ gives
\[
  x_B+\bar A_Nx_N=\bar b,\qquad
  \bar A=A_B^{-1}A,\qquad\bar b=A_B^{-1}b.
\]
These row operations preserve every solution of the equality system. They
change its representation, not its decisions. Substitute the basic equations
into the maximization objective:
\[
  z=c^Tx=z_0-q_N^Tx_N,\qquad
  z_0=c_B^T\bar b,\qquad q_j=c_B^T\bar a_j-c_j.
\]
Thus $q_j<0$ means that increasing nonbasic $x_j$ increases the objective.
The alternative convention $c_j-z_j$ is $-q_j$ and reverses this test.
The bottom row of each tableau lists $z_j-c_j$; its RHS cell records the current
objective $z_0$, rather than another equation coefficient. The $c_B$ column gives the objective coefficient of each basic
variable. The $z_j$ row sums each constraint column with these weights;
subtract the displayed $c_j$ row to obtain $q_j=z_j-c_j$. The same
weighted sum on the RHS gives $z_0$. In Phase I these are the auxiliary
objective coefficients; in Phase II they are the original ones. Basic indicators
are always zero and do not by themselves imply alternative optima.

When $x_k$ enters at value $\theta\geq0$, the basic variables become
$x_B(\theta)=\bar b-\theta\bar a_k$. A row with $\bar a_{ik}>0$
limits the step to $\bar b_i/\bar a_{ik}$, because its basic variable decreases.
A row with a zero or negative coefficient supplies no upper bound. Ratios are calculated in the column to the right of the RHS;
$\longrightarrow$ in the separate Out column marks the chosen leaving row.
The indicator row is labelled $z_j-c_j$; $q_j$ abbreviates this difference
in formulas. Vertical rules separate $c_B$, the basic-variable labels,
the coefficient columns, and the RHS. Variable-column headings keep their
order across pivots; only the basic-variable labels change. Artificial
columns are removed after Phase I without reordering the retained columns.
The minimum of the positive-coefficient ratios makes one basic variable zero
and identifies a leaving variable. We choose a most negative nonbasic
indicator, breaking entering ties by smaller variable index, for the improving
pivots below. Zero-indicator pivots are used afterwards to exhibit alternatives.

If the pivot is in row $r$, the complete augmented-row update is
\[
  R'_r=R_r/\bar a_{rk},\qquad
  R'_i=R_i-\bar a_{ik}R'_r\quad(i\ne r).
\]
This changes the entering column into an identity column. Recomputing $q$
with the new basic costs gives $z'_0=z_0-q_k\theta$.
If $q_k<0$ but no row coefficient is positive, there is no leaving variable:
$\theta$ can grow without bound. The full direction is
$d_k=1$, $d_B=-\bar a_k$, with all other nonbasic components zero.
It satisfies $d\geq0$, $Ad=0$ and $c^Td=-q_k>0$.
Its first three coordinates give a growth ray in the original variables.

For a finite optimum, $q\geq0$ gives the global bound
$z=z_0-q_N^Tx_N\leq z_0$. To claim another optimizer from a zero indicator,
we still need a positive feasible step that changes the original point.
An explicit optimal segment supplies stronger evidence than merely reporting
a different basis.

\section{Problem 1(a): a unique finite maximum}\label{sec:a}
\paragraph{Task statement.}
\[
  \max(-x_1+x_2+3x_3),\qquad
  \begin{cases}
    -4x_1+2x_2-x_3\leq2, \\
    2x_1-x_2+x_3\leq1,   \\
    3x_1+x_3\leq5,       \\
    x_1,x_2,x_3\geq0.
  \end{cases}
\]
Add slacks $x_4,x_5,x_6$ to the three rows. The origin is feasible, and its
slack basis $B=(4,5,6)$ has basic values $(2,1,5)$.
\begin{center}
  \begin{tabular}{r@{\hspace{0.2em}}|@{\hspace{0.2em}}c@{\enspace}|@{\enspace}rrrrrr@{\quad}|@{\quad}r@{\qquad\quad}l@{\quad}c}\cmidrule(r{1.5em}){1-9}
    $c_B$ & Basic & $x_1$ & $x_2$ & $x_3$ & $x_4$ & $x_5$ & $x_6$ & RHS & Ratio ($x_3\uparrow$) & Out\\
    & $c_j$ & $-1$ & $1$ & $3$ & $0$ & $0$ & $0$ & & & \\
    \cmidrule(r{1.5em}){1-9}
    $0$ & $x_4$ & $-4$ & $2$ & $-1$ & $1$ & $0$ & $0$ & $2$ & $\text{--}$ & \\
    $0$ & $x_5$ & $2$ & $-1$ & $1$ & $0$ & $1$ & $0$ & $1$ & $1 / 1=1$ & $\longrightarrow$\\
    $0$ & $x_6$ & $3$ & $0$ & $1$ & $0$ & $0$ & $1$ & $5$ & $5 / 1=5$ & \\
    \cmidrule(r{1.5em}){1-9}
    & $z_j$ & $0$ & $0$ & $0$ & $0$ & $0$ & $0$ & $0$ & & \\
    & $z_j-c_j$ & $1$ & $-1$ & $-3$ & $0$ & $0$ & $0$ & $z_0=0$ & & \\
    \cmidrule(r{1.5em}){1-9}
  \end{tabular}
\end{center}
Enter $x_3$, since $q_3=-3$, and remove $x_5$ at $\theta=1$. Keep row 2, add it to row 1, and subtract
it from row 3. The objective increases by $3\cdot1$.
\begin{center}
  \begin{tabular}{r@{\hspace{0.2em}}|@{\hspace{0.2em}}c@{\enspace}|@{\enspace}rrrrrr@{\quad}|@{\quad}r@{\qquad\quad}l@{\quad}c}\cmidrule(r{1.5em}){1-9}
    $c_B$ & Basic & $x_1$ & $x_2$ & $x_3$ & $x_4$ & $x_5$ & $x_6$ & RHS & Ratio ($x_2\uparrow$) & Out\\
    & $c_j$ & $-1$ & $1$ & $3$ & $0$ & $0$ & $0$ & & & \\
    \cmidrule(r{1.5em}){1-9}
    $0$ & $x_4$ & $-2$ & $1$ & $0$ & $1$ & $1$ & $0$ & $3$ & $3 / 1=3$ & $\longrightarrow$\\
    $3$ & $x_3$ & $2$ & $-1$ & $1$ & $0$ & $1$ & $0$ & $1$ & $\text{--}$ & \\
    $0$ & $x_6$ & $1$ & $1$ & $0$ & $0$ & $-1$ & $1$ & $4$ & $4 / 1=4$ & \\
    \cmidrule(r{1.5em}){1-9}
    & $z_j$ & $6$ & $-3$ & $3$ & $0$ & $3$ & $0$ & $3$ & & \\
    & $z_j-c_j$ & $7$ & $-4$ & $0$ & $0$ & $3$ & $0$ & $z_0=3$ & & \\
    \cmidrule(r{1.5em}){1-9}
  \end{tabular}
\end{center}
Now enter $x_2$, since $q_2=-4$. With $x_4$ leaving at step 3, add row 1 to row 2 and subtract row 1
from row 3. The original point becomes $(0,3,4)$ and the value becomes 15.
\begin{center}
  \begin{tabular}{r@{\hspace{0.2em}}|@{\hspace{0.2em}}c@{\enspace}|@{\enspace}rrrrrr@{\quad}|@{\quad}r@{\qquad\quad}l@{\quad}c}\cmidrule(r{1.5em}){1-9}
    $c_B$ & Basic & $x_1$ & $x_2$ & $x_3$ & $x_4$ & $x_5$ & $x_6$ & RHS & Ratio ($x_1\uparrow$) & Out\\
    & $c_j$ & $-1$ & $1$ & $3$ & $0$ & $0$ & $0$ & & & \\
    \cmidrule(r{1.5em}){1-9}
    $1$ & $x_2$ & $-2$ & $1$ & $0$ & $1$ & $1$ & $0$ & $3$ & $\text{--}$ & \\
    $3$ & $x_3$ & $0$ & $0$ & $1$ & $1$ & $2$ & $0$ & $4$ & $\text{--}$ & \\
    $0$ & $x_6$ & $3$ & $0$ & $0$ & $-1$ & $-2$ & $1$ & $1$ & $1 / 3$ & $\longrightarrow$\\
    \cmidrule(r{1.5em}){1-9}
    & $z_j$ & $-2$ & $1$ & $3$ & $4$ & $7$ & $0$ & $15$ & & \\
    & $z_j-c_j$ & $-1$ & $0$ & $0$ & $4$ & $7$ & $0$ & $z_0=15$ & & \\
    \cmidrule(r{1.5em}){1-9}
  \end{tabular}
\end{center}
Enter $x_1$, whose indicator is $-1$, and remove $x_6$ at $\theta=1/3$.
Divide row 3 by 3 and add twice the normalized row to row 1.
\begin{center}
  \begin{tabular}{r@{\hspace{0.2em}}|@{\hspace{0.2em}}c@{\enspace}|@{\enspace}rrrrrr@{\quad}|@{\quad}r}\toprule
    $c_B$ & Basic & $x_1$ & $x_2$ & $x_3$ & $x_4$ & $x_5$ & $x_6$ & RHS\\
    & $c_j$ & $-1$ & $1$ & $3$ & $0$ & $0$ & $0$ & \\
    \midrule
    $1$ & $x_2$ & $0$ & $1$ & $0$ & $\tfrac{1}{3}$ & $-\tfrac{1}{3}$ & $\tfrac{2}{3}$ & $\tfrac{11}{3}$\\
    $3$ & $x_3$ & $0$ & $0$ & $1$ & $1$ & $2$ & $0$ & $4$\\
    $-1$ & $x_1$ & $1$ & $0$ & $0$ & $-\tfrac{1}{3}$ & $-\tfrac{2}{3}$ & $\tfrac{1}{3}$ & $\tfrac{1}{3}$\\
    \midrule
    & $z_j$ & $-1$ & $1$ & $3$ & $\tfrac{11}{3}$ & $\tfrac{19}{3}$ & $\tfrac{1}{3}$ & $\tfrac{46}{3}$\\
    & $z_j-c_j$ & $0$ & $0$ & $0$ & $\tfrac{11}{3}$ & $\tfrac{19}{3}$ & $\tfrac{1}{3}$ & $z_0=\tfrac{46}{3}$\\
    \bottomrule
  \end{tabular}
\end{center}
All indicators are nonnegative. The original optimizer and value are
\[
  x^*=\left(\tfrac13,\tfrac{11}{3},4\right),\qquad z^*=\tfrac{46}{3}.
\]
The original slacks vanish, so all three resource rows are tight.
The final objective identity holds for every feasible point:
\[
  -x_1+x_2+3x_3
  =\tfrac{46}{3}-\tfrac{11}{3}x_4-\tfrac{19}{3}x_5-\tfrac13x_6
  \leq\tfrac{46}{3}.
\]
Equality forces $x_4=x_5=x_6=0$. Solving the three tight original rows
yields the displayed point uniquely. There are no alternative optimizers.

\section{Problem 1(b): an unbounded feasible set with a finite maximum}\label{sec:b}
\paragraph{Task statement.}
\[
  \max(2x_1-x_2+3x_3),\qquad
  \begin{cases}
    -3x_1-x_2-2x_3\leq3, \\
    3x_1+x_3\leq2,       \\
    -x_2+x_3\leq1,       \\
    x_1,x_2,x_3\geq0.
  \end{cases}
\]
Add slacks $x_4,x_5,x_6$. The first inequality is automatically satisfied
by nonnegative original variables, but retaining it gives a harmless basic
slack. Start at the origin with $B=(4,5,6)$.
\begin{center}
  \begin{tabular}{r@{\hspace{0.2em}}|@{\hspace{0.2em}}c@{\enspace}|@{\enspace}rrrrrr@{\quad}|@{\quad}r@{\qquad\quad}l@{\quad}c}\cmidrule(r{1.5em}){1-9}
    $c_B$ & Basic & $x_1$ & $x_2$ & $x_3$ & $x_4$ & $x_5$ & $x_6$ & RHS & Ratio ($x_3\uparrow$) & Out\\
    & $c_j$ & $2$ & $-1$ & $3$ & $0$ & $0$ & $0$ & & & \\
    \cmidrule(r{1.5em}){1-9}
    $0$ & $x_4$ & $-3$ & $-1$ & $-2$ & $1$ & $0$ & $0$ & $3$ & $\text{--}$ & \\
    $0$ & $x_5$ & $3$ & $0$ & $1$ & $0$ & $1$ & $0$ & $2$ & $2 / 1=2$ & \\
    $0$ & $x_6$ & $0$ & $-1$ & $1$ & $0$ & $0$ & $1$ & $1$ & $1 / 1=1$ & $\longrightarrow$\\
    \cmidrule(r{1.5em}){1-9}
    & $z_j$ & $0$ & $0$ & $0$ & $0$ & $0$ & $0$ & $0$ & & \\
    & $z_j-c_j$ & $-2$ & $1$ & $-3$ & $0$ & $0$ & $0$ & $z_0=0$ & & \\
    \cmidrule(r{1.5em}){1-9}
  \end{tabular}
\end{center}
Enter $x_3$ and remove $x_6$ with step 1. Add twice row 3 to row 1
and subtract row 3 from row 2.
\begin{center}
  \begin{tabular}{r@{\hspace{0.2em}}|@{\hspace{0.2em}}c@{\enspace}|@{\enspace}rrrrrr@{\quad}|@{\quad}r@{\qquad\quad}l@{\quad}c}\cmidrule(r{1.5em}){1-9}
    $c_B$ & Basic & $x_1$ & $x_2$ & $x_3$ & $x_4$ & $x_5$ & $x_6$ & RHS & Ratio ($x_1\uparrow$) & Out\\
    & $c_j$ & $2$ & $-1$ & $3$ & $0$ & $0$ & $0$ & & & \\
    \cmidrule(r{1.5em}){1-9}
    $0$ & $x_4$ & $-3$ & $-3$ & $0$ & $1$ & $0$ & $2$ & $5$ & $\text{--}$ & \\
    $0$ & $x_5$ & $3$ & $1$ & $0$ & $0$ & $1$ & $-1$ & $1$ & $1 / 3$ & $\longrightarrow$\\
    $3$ & $x_3$ & $0$ & $-1$ & $1$ & $0$ & $0$ & $1$ & $1$ & $\text{--}$ & \\
    \cmidrule(r{1.5em}){1-9}
    & $z_j$ & $0$ & $-3$ & $3$ & $0$ & $0$ & $3$ & $3$ & & \\
    & $z_j-c_j$ & $-2$ & $-2$ & $0$ & $0$ & $0$ & $3$ & $z_0=3$ & & \\
    \cmidrule(r{1.5em}){1-9}
  \end{tabular}
\end{center}
The indicators $q_1=q_2=-2$ tie; choose $x_1$ by smaller index.
The step is $1/3$.
Divide row 2 by 3 and add three copies to row 1; $x_5$ leaves.
\begin{center}
  \begin{tabular}{r@{\hspace{0.2em}}|@{\hspace{0.2em}}c@{\enspace}|@{\enspace}rrrrrr@{\quad}|@{\quad}r@{\qquad\quad}l@{\quad}c}\cmidrule(r{1.5em}){1-9}
    $c_B$ & Basic & $x_1$ & $x_2$ & $x_3$ & $x_4$ & $x_5$ & $x_6$ & RHS & Ratio ($x_2\uparrow$) & Out\\
    & $c_j$ & $2$ & $-1$ & $3$ & $0$ & $0$ & $0$ & & & \\
    \cmidrule(r{1.5em}){1-9}
    $0$ & $x_4$ & $0$ & $-2$ & $0$ & $1$ & $1$ & $1$ & $6$ & $\text{--}$ & \\
    $2$ & $x_1$ & $1$ & $\tfrac{1}{3}$ & $0$ & $0$ & $\tfrac{1}{3}$ & $-\tfrac{1}{3}$ & $\tfrac{1}{3}$ & $(1 / 3) / (1 / 3)=1$ & $\longrightarrow$\\
    $3$ & $x_3$ & $0$ & $-1$ & $1$ & $0$ & $0$ & $1$ & $1$ & $\text{--}$ & \\
    \cmidrule(r{1.5em}){1-9}
    & $z_j$ & $2$ & $-\tfrac{7}{3}$ & $3$ & $0$ & $\tfrac{2}{3}$ & $\tfrac{7}{3}$ & $\tfrac{11}{3}$ & & \\
    & $z_j-c_j$ & $0$ & $-\tfrac{4}{3}$ & $0$ & $0$ & $\tfrac{2}{3}$ & $\tfrac{7}{3}$ & $z_0=\tfrac{11}{3}$ & & \\
    \cmidrule(r{1.5em}){1-9}
  \end{tabular}
\end{center}
Enter $x_2$, since $q_2=-4/3$, and remove $x_1$ with step 1.
Multiply row 2 by 3,
add twice that normalized row to row 1, and add it to row 3.
\begin{center}
  \begin{tabular}{r@{\hspace{0.2em}}|@{\hspace{0.2em}}c@{\enspace}|@{\enspace}rrrrrr@{\quad}|@{\quad}r}\toprule
    $c_B$ & Basic & $x_1$ & $x_2$ & $x_3$ & $x_4$ & $x_5$ & $x_6$ & RHS\\
    & $c_j$ & $2$ & $-1$ & $3$ & $0$ & $0$ & $0$ & \\
    \midrule
    $0$ & $x_4$ & $6$ & $0$ & $0$ & $1$ & $3$ & $-1$ & $8$\\
    $-1$ & $x_2$ & $3$ & $1$ & $0$ & $0$ & $1$ & $-1$ & $1$\\
    $3$ & $x_3$ & $3$ & $0$ & $1$ & $0$ & $1$ & $0$ & $2$\\
    \midrule
    & $z_j$ & $6$ & $-1$ & $3$ & $0$ & $2$ & $1$ & $5$\\
    & $z_j-c_j$ & $4$ & $0$ & $0$ & $0$ & $2$ & $1$ & $z_0=5$\\
    \bottomrule
  \end{tabular}
\end{center}
The point is $x^*=(0,1,2)$, with maximum $z^*=5$ and slacks $(8,0,0)$.
For all feasible points the final objective is
\[
  2x_1-x_2+3x_3=5-4x_1-2x_5-x_6\leq5.
\]
Equality forces $x_1=x_5=x_6=0$. The second row then fixes $x_3=2$,
and the third fixes $x_2=1$; the optimizer is unique.
The zero indicator of basic $x_4$ supplies no alternative decision.

The feasible set itself is unbounded: $(0,t,0)$ is feasible for every $t\geq0$.
Its objective is $-t$, so that ray cannot improve a maximization objective.
Unboundedness of the feasible set does not imply an unbounded maximum.

\section{Problem 1(c): an entire segment of optimal solutions}\label{sec:c}
\paragraph{Task statement.}
\[
  \max(-4x_1+2x_2+3x_3),\qquad
  \begin{cases}
    2x_1-x_2+x_3\leq2,    \\
    -6x_1+3x_2-2x_3\leq2, \\
    -7x_1+4x_2-3x_3\leq2, \\
    x_1,x_2,x_3\geq0.
  \end{cases}
\]
Add slacks $x_4,x_5,x_6$ and start with $B=(4,5,6)$ at the origin.
\begin{center}
  \begin{tabular}{r@{\hspace{0.2em}}|@{\hspace{0.2em}}c@{\enspace}|@{\enspace}rrrrrr@{\quad}|@{\quad}r@{\qquad\quad}l@{\quad}c}\cmidrule(r{1.5em}){1-9}
    $c_B$ & Basic & $x_1$ & $x_2$ & $x_3$ & $x_4$ & $x_5$ & $x_6$ & RHS & Ratio ($x_3\uparrow$) & Out\\
    & $c_j$ & $-4$ & $2$ & $3$ & $0$ & $0$ & $0$ & & & \\
    \cmidrule(r{1.5em}){1-9}
    $0$ & $x_4$ & $2$ & $-1$ & $1$ & $1$ & $0$ & $0$ & $2$ & $2 / 1=2$ & $\longrightarrow$\\
    $0$ & $x_5$ & $-6$ & $3$ & $-2$ & $0$ & $1$ & $0$ & $2$ & $\text{--}$ & \\
    $0$ & $x_6$ & $-7$ & $4$ & $-3$ & $0$ & $0$ & $1$ & $2$ & $\text{--}$ & \\
    \cmidrule(r{1.5em}){1-9}
    & $z_j$ & $0$ & $0$ & $0$ & $0$ & $0$ & $0$ & $0$ & & \\
    & $z_j-c_j$ & $4$ & $-2$ & $-3$ & $0$ & $0$ & $0$ & $z_0=0$ & & \\
    \cmidrule(r{1.5em}){1-9}
  \end{tabular}
\end{center}
Enter $x_3$ and remove $x_4$ at $\theta=2$. Add twice row 1 to row 2 and three times row 1 to row 3.
\begin{center}
  \begin{tabular}{r@{\hspace{0.2em}}|@{\hspace{0.2em}}c@{\enspace}|@{\enspace}rrrrrr@{\quad}|@{\quad}r@{\qquad\quad}l@{\quad}c}\cmidrule(r{1.5em}){1-9}
    $c_B$ & Basic & $x_1$ & $x_2$ & $x_3$ & $x_4$ & $x_5$ & $x_6$ & RHS & Ratio ($x_2\uparrow$) & Out\\
    & $c_j$ & $-4$ & $2$ & $3$ & $0$ & $0$ & $0$ & & & \\
    \cmidrule(r{1.5em}){1-9}
    $3$ & $x_3$ & $2$ & $-1$ & $1$ & $1$ & $0$ & $0$ & $2$ & $\text{--}$ & \\
    $0$ & $x_5$ & $-2$ & $1$ & $0$ & $2$ & $1$ & $0$ & $6$ & $6 / 1=6$ & $\longrightarrow$\\
    $0$ & $x_6$ & $-1$ & $1$ & $0$ & $3$ & $0$ & $1$ & $8$ & $8 / 1=8$ & \\
    \cmidrule(r{1.5em}){1-9}
    & $z_j$ & $6$ & $-3$ & $3$ & $3$ & $0$ & $0$ & $6$ & & \\
    & $z_j-c_j$ & $10$ & $-5$ & $0$ & $3$ & $0$ & $0$ & $z_0=6$ & & \\
    \cmidrule(r{1.5em}){1-9}
  \end{tabular}
\end{center}
Enter $x_2$, since $q_2=-5$, and remove $x_5$ with step 6.
Add row 2 to row 1 and subtract row 2 from row 3.
\begin{center}
  \begin{tabular}{r@{\hspace{0.2em}}|@{\hspace{0.2em}}c@{\enspace}|@{\enspace}rrrrrr@{\quad}|@{\quad}r@{\qquad\quad}l@{\quad}c}\cmidrule(r{1.5em}){1-9}
    $c_B$ & Basic & $x_1$ & $x_2$ & $x_3$ & $x_4$ & $x_5$ & $x_6$ & RHS & Ratio ($x_1\uparrow$) & Out\\
    & $c_j$ & $-4$ & $2$ & $3$ & $0$ & $0$ & $0$ & & & \\
    \cmidrule(r{1.5em}){1-9}
    $3$ & $x_3$ & $0$ & $0$ & $1$ & $3$ & $1$ & $0$ & $8$ & $\text{--}$ & \\
    $2$ & $x_2$ & $-2$ & $1$ & $0$ & $2$ & $1$ & $0$ & $6$ & $\text{--}$ & \\
    $0$ & $x_6$ & $1$ & $0$ & $0$ & $1$ & $-1$ & $1$ & $2$ & $2 / 1=2$ & $\longrightarrow$\\
    \cmidrule(r{1.5em}){1-9}
    & $z_j$ & $-4$ & $2$ & $3$ & $13$ & $5$ & $0$ & $36$ & & \\
    & $z_j-c_j$ & $0$ & $0$ & $0$ & $13$ & $5$ & $0$ & $z_0=36$ & & \\
    \cmidrule(r{1.5em}){1-9}
  \end{tabular}
\end{center}
All indicators are nonnegative, so $(0,6,8)$ is optimal with value 36.
Here $x_1$ is \emph{nonbasic} with $q_1=0$. To test for an actual alternative,
let it enter with step $\theta=2>0$, removing $x_6$.
The original point changes and the objective change is $-q_1\theta=0$.
Add twice row 3 to row 2; $x_6$ leaves.
\begin{center}
  \begin{tabular}{r@{\hspace{0.2em}}|@{\hspace{0.2em}}c@{\enspace}|@{\enspace}rrrrrr@{\quad}|@{\quad}r}\toprule
    $c_B$ & Basic & $x_1$ & $x_2$ & $x_3$ & $x_4$ & $x_5$ & $x_6$ & RHS\\
    & $c_j$ & $-4$ & $2$ & $3$ & $0$ & $0$ & $0$ & \\
    \midrule
    $3$ & $x_3$ & $0$ & $0$ & $1$ & $3$ & $1$ & $0$ & $8$\\
    $2$ & $x_2$ & $0$ & $1$ & $0$ & $4$ & $-1$ & $2$ & $10$\\
    $-4$ & $x_1$ & $1$ & $0$ & $0$ & $1$ & $-1$ & $1$ & $2$\\
    \midrule
    & $z_j$ & $-4$ & $2$ & $3$ & $13$ & $5$ & $0$ & $36$\\
    & $z_j-c_j$ & $0$ & $0$ & $0$ & $13$ & $5$ & $0$ & $z_0=36$\\
    \bottomrule
  \end{tabular}
\end{center}
The second endpoint is $(2,10,8)$, also of value 36.
The global certificate and complete set of optimizers are
\[
  -4x_1+2x_2+3x_3=36-13x_4-5x_5\leq36,
\]
\begin{equation}\label{eq:c-segment}
  x(t)=(t,6+2t,8),\qquad0\leq t\leq2.
\end{equation}
Indeed, equality forces $x_4=x_5=0$. Their tight original rows give
$x_3=8$ and $x_2=6+2x_1$. The third slack is $x_6=2-x_1$,
so nonnegativity gives exactly $0\leq x_1\leq2$.
This derives every optimizer, including nonbasic interior points of the segment.

\section{Problem 1(d): starting at the supplied point}\label{sec:d}
\paragraph{Task statement.}
\[
  \max(x_1+2x_2+x_3),\qquad
  \begin{cases}
    2x_2+x_3\geq4,      \\
    x_2+x_3\leq10,      \\
    x_1+2x_2+2x_3\leq8, \\
    x_1,x_2,x_3\geq0,
  \end{cases}
  \qquad x_0=(0,2,0).
\]
For the first row subtract the surplus $x_4=2x_2+x_3-4\geq0$;
for the other rows add slacks $x_5,x_6$. The equality equations are
\[
  2x_2+x_3-x_4=4,\qquad x_2+x_3+x_5=10,\qquad
  x_1+2x_2+2x_3+x_6=8.
\]
The supplied point is feasible and its auxiliary values are $(0,8,4)$.
Choose $B=(2,5,6)$: the columns are independent, and the positive basic
values are $(2,8,4)$. No artificial variable is needed merely because a
$\geq$ row occurs. Divide the first equation by 2, subtract the resulting
row from row 2, and subtract twice it from row 3.
\begin{center}
  \begin{tabular}{r@{\hspace{0.2em}}|@{\hspace{0.2em}}c@{\enspace}|@{\enspace}rrrrrr@{\quad}|@{\quad}r@{\qquad\quad}l@{\quad}c}\cmidrule(r{1.5em}){1-9}
    $c_B$ & Basic & $x_1$ & $x_2$ & $x_3$ & $x_4$ & $x_5$ & $x_6$ & RHS & Ratio ($x_1\uparrow$) & Out\\
    & $c_j$ & $1$ & $2$ & $1$ & $0$ & $0$ & $0$ & & & \\
    \cmidrule(r{1.5em}){1-9}
    $2$ & $x_2$ & $0$ & $1$ & $\tfrac{1}{2}$ & $-\tfrac{1}{2}$ & $0$ & $0$ & $2$ & $\text{--}$ & \\
    $0$ & $x_5$ & $0$ & $0$ & $\tfrac{1}{2}$ & $\tfrac{1}{2}$ & $1$ & $0$ & $8$ & $\text{--}$ & \\
    $0$ & $x_6$ & $1$ & $0$ & $1$ & $1$ & $0$ & $1$ & $4$ & $4 / 1=4$ & $\longrightarrow$\\
    \cmidrule(r{1.5em}){1-9}
    & $z_j$ & $0$ & $2$ & $1$ & $-1$ & $0$ & $0$ & $4$ & & \\
    & $z_j-c_j$ & $-1$ & $0$ & $0$ & $-1$ & $0$ & $0$ & $z_0=4$ & & \\
    \cmidrule(r{1.5em}){1-9}
  \end{tabular}
\end{center}
The indicators of $x_1$ and surplus $x_4$ tie at $-1$; choose $x_1$.
Remove $x_6$ at $\theta=4$. That row already contains the entering identity column,
so the constraint coefficients do not change; the basic label and objective do.
\begin{center}
  \begin{tabular}{r@{\hspace{0.2em}}|@{\hspace{0.2em}}c@{\enspace}|@{\enspace}rrrrrr@{\quad}|@{\quad}r@{\qquad\quad}l@{\quad}c}\cmidrule(r{1.5em}){1-9}
    $c_B$ & Basic & $x_1$ & $x_2$ & $x_3$ & $x_4$ & $x_5$ & $x_6$ & RHS & Ratio ($x_4\uparrow$) & Out\\
    & $c_j$ & $1$ & $2$ & $1$ & $0$ & $0$ & $0$ & & & \\
    \cmidrule(r{1.5em}){1-9}
    $2$ & $x_2$ & $0$ & $1$ & $\tfrac{1}{2}$ & $-\tfrac{1}{2}$ & $0$ & $0$ & $2$ & $\text{--}$ & \\
    $0$ & $x_5$ & $0$ & $0$ & $\tfrac{1}{2}$ & $\tfrac{1}{2}$ & $1$ & $0$ & $8$ & $8 / (1 / 2)=16$ & \\
    $1$ & $x_1$ & $1$ & $0$ & $1$ & $1$ & $0$ & $1$ & $4$ & $4 / 1=4$ & $\longrightarrow$\\
    \cmidrule(r{1.5em}){1-9}
    & $z_j$ & $1$ & $2$ & $2$ & $0$ & $0$ & $1$ & $8$ & & \\
    & $z_j-c_j$ & $0$ & $0$ & $1$ & $0$ & $0$ & $1$ & $z_0=8$ & & \\
    \cmidrule(r{1.5em}){1-9}
  \end{tabular}
\end{center}
The original point $(4,2,0)$ is optimal with value 8. Its objective identity is
\[
  x_1+2x_2+x_3=8-x_3-x_6\leq8.
\]
To find an alternative, enter nonbasic $x_4$ with $q_4=0$.
Remove $x_1$ with step 4. Add half row 3 to row 1 and subtract half row 3 from row 2.
\begin{center}
  \begin{tabular}{r@{\hspace{0.2em}}|@{\hspace{0.2em}}c@{\enspace}|@{\enspace}rrrrrr@{\quad}|@{\quad}r}\toprule
    $c_B$ & Basic & $x_1$ & $x_2$ & $x_3$ & $x_4$ & $x_5$ & $x_6$ & RHS\\
    & $c_j$ & $1$ & $2$ & $1$ & $0$ & $0$ & $0$ & \\
    \midrule
    $2$ & $x_2$ & $\tfrac{1}{2}$ & $1$ & $1$ & $0$ & $0$ & $\tfrac{1}{2}$ & $4$\\
    $0$ & $x_5$ & $-\tfrac{1}{2}$ & $0$ & $0$ & $0$ & $1$ & $-\tfrac{1}{2}$ & $6$\\
    $0$ & $x_4$ & $1$ & $0$ & $1$ & $1$ & $0$ & $1$ & $4$\\
    \midrule
    & $z_j$ & $1$ & $2$ & $2$ & $0$ & $0$ & $1$ & $8$\\
    & $z_j-c_j$ & $0$ & $0$ & $1$ & $0$ & $0$ & $1$ & $z_0=8$\\
    \bottomrule
  \end{tabular}
\end{center}
The other endpoint is $(0,4,0)$, still of value 8. For any optimum,
the certificate forces $x_3=x_6=0$, so $x_1+2x_2=8$.
The first row requires $x_2\geq2$ and $x_1\geq0$ requires $x_2\leq4$.
Therefore the complete optimal segment is
\begin{equation}\label{eq:d-segment}
  x(t)=(8-2t,t,0),\qquad2\leq t\leq4.
\end{equation}
The second row remains slack throughout, since $x_5=10-t>0$.

\section{Problem 1(e): reject the supplied start and exhibit a growth ray}\label{sec:e}
\paragraph{Task statement.}
\[
  \max(2x_1+x_2+x_3),\qquad
  \begin{cases}
    2x_1+3x_2-x_3\leq9, \\
    2x_2+x_3\geq4,      \\
    x_1=6,              \\
    x_1,x_2,x_3\geq0,
  \end{cases}
  \qquad x_0=(0,0,6).
\]
The supplied $x_0$ violates the equality $x_1=6$, although it satisfies the
two inequalities and nonnegativity. It cannot start primal simplex. Keep the
equality and find a feasible point for the stated problem: with $x_1=6$
and $x_2=0$, the inequalities become $x_3\geq3$ and $x_3\geq4$.
Choose $x_3=4$.

Add first-row slack $x_4$ and subtract second-row surplus $x_5$:
\[
  2x_1+3x_2-x_3+x_4=9,\qquad2x_2+x_3-x_5=4,\qquad x_1=6.
\]
The feasible augmented point is $(6,0,4,1,0)$.
The independent columns $B=(4,3,1)$ give basic values $(1,4,6)$.
Subtract twice the equality row from row 1, then add row 2 to row 1:
\begin{center}
  \begin{tabular}{r@{\hspace{0.2em}}|@{\hspace{0.2em}}c@{\enspace}|@{\enspace}rrrrr@{\quad}|@{\quad}r@{\qquad\quad}l@{\quad}c}\cmidrule(r{1.5em}){1-8}
    $c_B$ & Basic & $x_1$ & $x_2$ & $x_3$ & $x_4$ & $x_5$ & RHS & Ratio ($x_5\uparrow$) & Out\\
    & $c_j$ & $2$ & $1$ & $1$ & $0$ & $0$ & & & \\
    \cmidrule(r{1.5em}){1-8}
    $0$ & $x_4$ & $0$ & $5$ & $0$ & $1$ & $-1$ & $1$ & $\text{--}$ & \\
    $1$ & $x_3$ & $0$ & $2$ & $1$ & $0$ & $-1$ & $4$ & $\text{--}$ & \\
    $2$ & $x_1$ & $1$ & $0$ & $0$ & $0$ & $0$ & $6$ & $\text{--}$ & \\
    \cmidrule(r{1.5em}){1-8}
    & $z_j$ & $2$ & $2$ & $1$ & $0$ & $-1$ & $16$ & & \\
    & $z_j-c_j$ & $0$ & $1$ & $0$ & $0$ & $-1$ & $z_0=16$ & & \\
    \cmidrule(r{1.5em}){1-8}
  \end{tabular}
\end{center}
Only $q_5=-1$ is negative. Its column is $(-1,-1,0)^T$,
so increasing $x_5$ \emph{increases} the first two basic values and leaves
$x_1$ unchanged. There is no positive coefficient and no leaving variable.
The augmented direction and original growth ray are
\[
  d=(0,0,1,1,1),\qquad
  x(\theta)=(6,0,4)+\theta(0,0,1),\qquad\theta\geq0.
\]
For every $\theta\geq0$, the original left sides are
$8-\theta\leq9$, $4+\theta\geq4$ and $x_1=6$.
The objective is $16+\theta\to+\infty$.
Thus the problem is feasible and unbounded above; it has no finite optimum
whose alternatives could be classified. Rejection of the supplied point
is different from infeasibility of the problem.

\section{Problem 1(f): two supplied starts, one unbounded problem}\label{sec:f}
\paragraph{Task statement.}
\[
  \max(6x_1+2x_2-3x_3),\qquad
  \begin{cases}
    x_1+2x_2-3x_3\leq3,  \\
    2x_1+5x_2-5x_3\leq7, \\
    2x_1-3x_2-7x_3\leq8, \\
    x_1,x_2,x_3\geq0.
  \end{cases}
\]
\begin{otherlanguage}{polish}
  i) $x_0=(0,1,0)$, ii) $x_0=(3,0,0)$,
\end{otherlanguage}
Add slacks $x_4,x_5,x_6$ to the three rows. The different supplied points
require different initialization, even though the LP itself is identical.

\subsection{Start i: a feasible point that is not basic}\label{sec:f-i}
At $(0,1,0)$ the original row values are $(2,5,-3)$ and the slacks are
$(1,2,11)$. All constraints are satisfied. However, $x_2,x_4,x_5,x_6$
are positive: four positive augmented coordinates cannot fit into a basis of
three columns, whose nonbasic coordinates must all be zero. No choice of
basis makes this point a basic solution in this slack formulation.

We can use the supplied point to reach a feasible vertex before performing
primal simplex. Keep $x_1=x_3=0$ and increase $x_2=1+\tau$.
The slacks become
\[
  x_4=1-2\tau,\qquad x_5=2-5\tau,\qquad x_6=11+3\tau.
\]
Feasibility requires $0\leq\tau\leq\min(1/2,2/5)=2/5$.
At the first boundary, $x_2=7/5$ and the slacks are $(1/5,0,61/5)$.
The objective rises from 2 to $14/5$. This preliminary move is a feasible
initialization from a nonbasic point, not a simplex pivot from an existing basis.
Now $B=(4,2,6)$ is independent and has three positive basic values.
Divide original row 2 by 5, subtract twice the normalized row from row 1,
and add three times it to row 3:
\begin{center}
  \begin{tabular}{r@{\hspace{0.2em}}|@{\hspace{0.2em}}c@{\enspace}|@{\enspace}rrrrrr@{\quad}|@{\quad}r@{\qquad\quad}l@{\quad}c}\cmidrule(r{1.5em}){1-9}
    $c_B$ & Basic & $x_1$ & $x_2$ & $x_3$ & $x_4$ & $x_5$ & $x_6$ & RHS & Ratio ($x_1\uparrow$) & Out\\
    & $c_j$ & $6$ & $2$ & $-3$ & $0$ & $0$ & $0$ & & & \\
    \cmidrule(r{1.5em}){1-9}
    $0$ & $x_4$ & $\tfrac{1}{5}$ & $0$ & $-1$ & $1$ & $-\tfrac{2}{5}$ & $0$ & $\tfrac{1}{5}$ & $(1 / 5) / (1 / 5)=1$ & $\longrightarrow$\\
    $2$ & $x_2$ & $\tfrac{2}{5}$ & $1$ & $-1$ & $0$ & $\tfrac{1}{5}$ & $0$ & $\tfrac{7}{5}$ & $(7 / 5) / (2 / 5)=7 / 2$ & \\
    $0$ & $x_6$ & $\tfrac{16}{5}$ & $0$ & $-10$ & $0$ & $\tfrac{3}{5}$ & $1$ & $\tfrac{61}{5}$ & $(61 / 5) / (16 / 5)=61 / 16$ & \\
    \cmidrule(r{1.5em}){1-9}
    & $z_j$ & $\tfrac{4}{5}$ & $2$ & $-2$ & $0$ & $\tfrac{2}{5}$ & $0$ & $\tfrac{14}{5}$ & & \\
    & $z_j-c_j$ & $-\tfrac{26}{5}$ & $0$ & $1$ & $0$ & $\tfrac{2}{5}$ & $0$ & $z_0=\tfrac{14}{5}$ & & \\
    \cmidrule(r{1.5em}){1-9}
  \end{tabular}
\end{center}
Enter $x_1$ with $q_1=-26/5$ and remove $x_4$ at $\theta=1$.
Multiply row 1 by 5, subtract $(2/5)$ times that normalized row from row 2,
and subtract $(16/5)$ times it from row 3.
\begin{center}
  \begin{tabular}{r@{\hspace{0.2em}}|@{\hspace{0.2em}}c@{\enspace}|@{\enspace}rrrrrr@{\quad}|@{\quad}r@{\qquad\quad}l@{\quad}c}\cmidrule(r{1.5em}){1-9}
    $c_B$ & Basic & $x_1$ & $x_2$ & $x_3$ & $x_4$ & $x_5$ & $x_6$ & RHS & Ratio ($x_3\uparrow$) & Out\\
    & $c_j$ & $6$ & $2$ & $-3$ & $0$ & $0$ & $0$ & & & \\
    \cmidrule(r{1.5em}){1-9}
    $6$ & $x_1$ & $1$ & $0$ & $-5$ & $5$ & $-2$ & $0$ & $1$ & $\text{--}$ & \\
    $2$ & $x_2$ & $0$ & $1$ & $1$ & $-2$ & $1$ & $0$ & $1$ & $1 / 1=1$ & $\longrightarrow$\\
    $0$ & $x_6$ & $0$ & $0$ & $6$ & $-16$ & $7$ & $1$ & $9$ & $9 / 6=3 / 2$ & \\
    \cmidrule(r{1.5em}){1-9}
    & $z_j$ & $6$ & $2$ & $-28$ & $26$ & $-10$ & $0$ & $8$ & & \\
    & $z_j-c_j$ & $0$ & $0$ & $-25$ & $26$ & $-10$ & $0$ & $z_0=8$ & & \\
    \cmidrule(r{1.5em}){1-9}
  \end{tabular}
\end{center}
The point is $(1,1,0)$ and the objective is 8. Enter $x_3$ with
$q_3=-25$, and remove $x_2$ with step 1.
Add five copies of row 2 to row 1 and subtract six copies from row 3:
\begin{center}
  \begin{tabular}{r@{\hspace{0.2em}}|@{\hspace{0.2em}}c@{\enspace}|@{\enspace}rrrrrr@{\quad}|@{\quad}r@{\qquad\quad}l@{\quad}c}\cmidrule(r{1.5em}){1-9}
    $c_B$ & Basic & $x_1$ & $x_2$ & $x_3$ & $x_4$ & $x_5$ & $x_6$ & RHS & Ratio ($x_4\uparrow$) & Out\\
    & $c_j$ & $6$ & $2$ & $-3$ & $0$ & $0$ & $0$ & & & \\
    \cmidrule(r{1.5em}){1-9}
    $6$ & $x_1$ & $1$ & $5$ & $0$ & $-5$ & $3$ & $0$ & $6$ & $\text{--}$ & \\
    $-3$ & $x_3$ & $0$ & $1$ & $1$ & $-2$ & $1$ & $0$ & $1$ & $\text{--}$ & \\
    $0$ & $x_6$ & $0$ & $-6$ & $0$ & $-4$ & $1$ & $1$ & $3$ & $\text{--}$ & \\
    \cmidrule(r{1.5em}){1-9}
    & $z_j$ & $6$ & $27$ & $-3$ & $-24$ & $15$ & $0$ & $33$ & & \\
    & $z_j-c_j$ & $0$ & $25$ & $0$ & $-24$ & $15$ & $0$ & $z_0=33$ & & \\
    \cmidrule(r{1.5em}){1-9}
  \end{tabular}
\end{center}
This gives $(6,0,1)$, value 33. The negative indicator $q_4=-24$
means this point is not optimal. The next entering variable is the slack
$x_4$, whose whole column is nonpositive. The growth ray is derived in
Section~\ref{sec:f-ray}.

\subsection{Start ii: the supplied point is already basic}\label{sec:f-ii}
At $(3,0,0)$ the slacks are $(0,1,2)$. Columns $B=(1,5,6)$ form an
independent basis with positive values $(3,1,2)$.
Keep original row 1, subtract twice it from rows 2 and 3, and recompute the
objective using the basic cost of $x_1$:
\begin{center}
  \begin{tabular}{r@{\hspace{0.2em}}|@{\hspace{0.2em}}c@{\enspace}|@{\enspace}rrrrrr@{\quad}|@{\quad}r@{\qquad\quad}l@{\quad}c}\cmidrule(r{1.5em}){1-9}
    $c_B$ & Basic & $x_1$ & $x_2$ & $x_3$ & $x_4$ & $x_5$ & $x_6$ & RHS & Ratio ($x_3\uparrow$) & Out\\
    & $c_j$ & $6$ & $2$ & $-3$ & $0$ & $0$ & $0$ & & & \\
    \cmidrule(r{1.5em}){1-9}
    $6$ & $x_1$ & $1$ & $2$ & $-3$ & $1$ & $0$ & $0$ & $3$ & $\text{--}$ & \\
    $0$ & $x_5$ & $0$ & $1$ & $1$ & $-2$ & $1$ & $0$ & $1$ & $1 / 1=1$ & $\longrightarrow$\\
    $0$ & $x_6$ & $0$ & $-7$ & $-1$ & $-2$ & $0$ & $1$ & $2$ & $\text{--}$ & \\
    \cmidrule(r{1.5em}){1-9}
    & $z_j$ & $6$ & $12$ & $-18$ & $6$ & $0$ & $0$ & $18$ & & \\
    & $z_j-c_j$ & $0$ & $10$ & $-15$ & $6$ & $0$ & $0$ & $z_0=18$ & & \\
    \cmidrule(r{1.5em}){1-9}
  \end{tabular}
\end{center}
Enter $x_3$, whose indicator is $-15$.
Remove $x_5$ at $\theta=1$.
Add three copies of row 2 to row 1 and add row 2 to row 3.
\begin{center}
  \begin{tabular}{r@{\hspace{0.2em}}|@{\hspace{0.2em}}c@{\enspace}|@{\enspace}rrrrrr@{\quad}|@{\quad}r@{\qquad\quad}l@{\quad}c}\cmidrule(r{1.5em}){1-9}
    $c_B$ & Basic & $x_1$ & $x_2$ & $x_3$ & $x_4$ & $x_5$ & $x_6$ & RHS & Ratio ($x_4\uparrow$) & Out\\
    & $c_j$ & $6$ & $2$ & $-3$ & $0$ & $0$ & $0$ & & & \\
    \cmidrule(r{1.5em}){1-9}
    $6$ & $x_1$ & $1$ & $5$ & $0$ & $-5$ & $3$ & $0$ & $6$ & $\text{--}$ & \\
    $-3$ & $x_3$ & $0$ & $1$ & $1$ & $-2$ & $1$ & $0$ & $1$ & $\text{--}$ & \\
    $0$ & $x_6$ & $0$ & $-6$ & $0$ & $-4$ & $1$ & $1$ & $3$ & $\text{--}$ & \\
    \cmidrule(r{1.5em}){1-9}
    & $z_j$ & $6$ & $27$ & $-3$ & $-24$ & $15$ & $0$ & $33$ & & \\
    & $z_j-c_j$ & $0$ & $25$ & $0$ & $-24$ & $15$ & $0$ & $z_0=33$ & & \\
    \cmidrule(r{1.5em}){1-9}
  \end{tabular}
\end{center}
The new basis is $(1,3,6)$, exactly the same canonical system reached
from start i. The original point is $(6,0,1)$ and its objective is 33.
The two different histories therefore lead to the same unboundedness test.

\subsection{A ray, not a last optimal point}\label{sec:f-ray}
At the shared final basis, $q_4=-24$ and
$\bar a_4=(-5,-2,-4)^T$. Increasing nonbasic slack $x_4$ by $\theta$
increases basic $(x_1,x_3,x_6)$ by $(5,2,4)\theta$.
No basic value decreases, so no finite ratio test bounds the step.
The full augmented direction and its original-variable projection are
\[
  d=(5,0,2,1,0,4),\qquad
  x(\theta)=(6,0,1)+\theta(5,0,2),\qquad\theta\geq0.
\]
The three original row values on this ray are
\[
  3-\theta\leq3,\qquad 7\leq7,\qquad5-4\theta\leq8.
\]
Nonnegativity holds for every $\theta\geq0$ and
\[
  6x_1(\theta)+2x_2(\theta)-3x_3(\theta)=33+24\theta\longrightarrow+\infty.
\]
The negative objective coefficient of $x_3$ does not prevent improvement:
its increase permits a sufficiently profitable increase in $x_1$.
A slack can enter a basis just like an original variable; projecting the
resulting direction explains the change in original decisions.
There is no finite optimal solution for either start.

\section{What the six outcomes teach}\label{sec:review}
\begin{center}
  \begin{tabular}{lll}\toprule
    Problem & Outcome & Original solution or family\\\midrule
    1(a) & Unique maximum $46/3$ & $(1/3,11/3,4)$\\
    1(b) & Unique maximum $5$ & $(0,1,2)$\\
    1(c) & Maximum $36$, alternatives & $(t,6+2t,8)$, $0\leq t\leq2$\\
    1(d) & Maximum $8$, alternatives & $(8-2t,t,0)$, $2\leq t\leq4$\\
    1(e) & Unbounded above & $(6,0,4)+\theta(0,0,1)$\\
    1(f) & Unbounded above & $(6,0,1)+\theta(5,0,2)$\\\bottomrule
  \end{tabular}
\end{center}
For the two rays, $\theta\geq0$. The growth rate depends on the scaling of
the direction; the existence of a positive rate is the unboundedness certificate.
Check domains and all original constraints before choosing a start. Feasibility
alone does not establish basicness. At a finite optimum use the indicator
identity to prove a global bound, then derive uniqueness or the entire optimal
face. For an unbounded problem provide both a feasible base point and a direction;
a failed ratio test without that interpretation is an incomplete explanation.

\section{Sources and editorial remarks}\label{sec:sources}
\noindent\textbf{Task source:} K. Junosza-Szaniawski,
\texttt{../execrices-2026-27/ProgMatc1.pdf}, PDF page 1, task 1(a)--(f).
The original Polish common instruction, mathematical data, supplied points and
variable domains are preserved. The stated equality $x_1=6$ in 1(e) is retained;
the supplied point fails that equality. The wording ``if possible'' permits
rejecting the start. No undocumented change to the problem is assumed.
Both source pages were inspected visually to confirm the statements.

\paragraph{Editorial additions and alignment.}
The hand derivations, explicit basis tests, full tableaus, complete optimal
segments and growth-ray certificates are added explanations. All tableaus and
pivot row operations were independently checked using exact rational arithmetic.
The convention $q=z-c$, slack/surplus meanings, and certificate arguments match
Chapters 2--4 of the reference notes. These are current supplied exercise tasks;
their presence establishes task coverage, not a claim about how much was taught
or explained in a particular class session.

\paragraph{MATLAB comparison.}
The actual companion is the primal-simplex MATLAB script:
\begin{center}
  \href{run:matlab/exercises/sheet_01_primal_simplex.m}{\texttt{matlab/exercises/sheet\_01\_primal\_simplex.m}}
\end{center}
Run it from the repository root:
\begin{lstlisting}
run('matlab/exercises/sheet_01_primal_simplex.m')
\end{lstlisting}
It uses base MATLAB, fixed supplied matrices and the prescribed pivot paths;
it does not call an optimization solver or select pivots automatically.
It recomputes each basis with backslash, checks every displayed tableau,
indicator, point, objective change and row operation, tests the supplied starts,
and verifies the finite bounds, optimal-segment endpoints and ray equations.
The absolute tolerance is $10^{-10}$ for these small dimensionless problems.
Exact hand fractions are compared with floating-point calculations on identical
inputs. Endpoint and direction identities prove the entire segments/rays;
additional sample points are numerical checks rather than substitutes for proof.

\textbf{Execution status: passed in MATLAB.}
The user supplied output on 2026-10-11 and confirmed version
26.2.0.3386108 (R2026b). All prescribed paths and certificate assertions passed.
The largest checked tableau/basis discrepancy was $7.11\cdot10^{-15}$,
below the absolute tolerance $10^{-10}$; this is an observed discrepancy for
these checks, not a universal solution-error bound.
Independent exact-arithmetic verification establishes the hand calculations.
The supplied output is retained in \path{docs/matlab-verification.md}.
\end{document}
